\chapter{Source Data and Audit Tables}\label{app:source}

This appendix reproduces the source simulations of Gilmore \& Stixrude \cite{gilmore2026} that
enter the audit and the recalibration, next to the predictions of the published fit. The
phase-coexistence runs contain 27 MgSiO$_3$ formula units and 92 to 219 hydrogen atoms; the
homogeneous runs contain 135 to 160 atoms. All
compositions are H$_2$ mole fractions $x$ on the oxide basis (Eq.~\ref{eq:x}).

\section{Phase-Coexistence Runs}

Table~\ref{tab:ed1} lists all 14 phase-coexistence runs. For runs that kept two phases it gives
the composition of the H-poor and H-rich phase from ED Table~1 and, where they differ, from the
digitized Figure~2a. For runs that became homogeneous it gives the bulk composition in
parentheses. The last column is the published critical temperature at the run pressure minus
the run temperature: positive values mean the model places the run below $T_C$.

\begin{table}[h]
\thesiscaption{Phase-coexistence runs of ED Table~1 and the published fit.}{Pressure
uncertainties in the last digits are given in parentheses. ``same'' means that the digitized
Figure~2a agrees with the table to within 0.001.}
\label{tab:ed1}
\small\centering
\begin{tabular}{@{}rrcllll@{}}
\toprule
$P$ (GPa) & $T$ (K) & Phases & ED Table~1 $x$ & Figure~2a $x$ & Published fit $x$ &
$T_C - T$ (K)\\
\midrule
\inputtable{tables/ed_table1.tex}
\bottomrule
\end{tabular}
\end{table}

In every two-phase run the per-phase atom numbers of ED Table~1 add up to the full simulation
cell, to within one atom. Each atom, including those in the two interfaces, is therefore
assigned to one of the two phases. For the 2.05 and 9.58\,GPa runs the table and the figure
agree to within \digMarker{}. For the 1.65 and 5.14\,GPa runs, which lie closest to the
critical curve, they differ by up to 0.07.

\section{Homogeneous Excess Properties}

Table~\ref{tab:homogeneous} lists the excess properties reconstructed from ED Table~2 per mole
of mixing units. The uncertainties include the shared endmember errors and the pressure-offset
systematic described in Section~\ref{sec:recalibration}. The enthalpy entries of different
compositions are correlated with coefficients of 0.48 to 0.76, the volumes with 0.48 to 0.78.

\begin{table}[h]
\thesiscaption{Excess properties of the homogeneous simulations at 5.06--5.08\,GPa and
4,000\,K}{(ED Table~2), with 1$\sigma$ uncertainties, next to the published fit.}
\label{tab:homogeneous}
\small\centering
\begin{tabular}{@{}lllllll@{}}
\toprule
& \multicolumn{2}{c}{$H_{\rm ex}$ (kJ\,mol$^{-1}$)} &
\multicolumn{2}{c}{$S_{\rm ex}$ (J\,mol$^{-1}$\,K$^{-1}$)} &
\multicolumn{2}{c}{$V_{\rm ex}$ (cm$^3$\,mol$^{-1}$)}\\
\cmidrule(lr){2-3}\cmidrule(lr){4-5}\cmidrule(lr){6-7}
$x$ & simulation & fit & simulation & fit & simulation & fit\\
\midrule
\inputtable{tables/homogeneous.tex}
\bottomrule
\end{tabular}
\end{table}
